%======================================================================
\chapter{Proposal}
%======================================================================

% The proposal is a document that outlines the topic of your project. Effectively, the proposal should describe why the
% scope of the project is appropriate for this course. The course instructor will provide feedback on the proposal
% following submission. Students are encouraged to be as detailed as possible in the proposal. The more details
% provided in the proposal, the more detailed feedback you will receive.

% The proposal should contain the following elements:
% 1. Background and objective(s) for the project.
% 2. Proposed method(s) from artificial intelligence.
% 3. Dataset or environment to be used.
% 4. Proposed validation/analysis strategy.
% 5. Description of the novelty of the project.
% 6. Schedule of milestones to be completed each week (beginning the week after the proposal is submitted).

% Page limit for proposals: 2 pages maximum
% Please use a standard page format (i.e. page size 8.5” x 11”, minimum ½” margins, minimum font size 10, no
% condensed fonts or spacing). Please submit the proposal as a PDF file.
% Proposals are to be submitted electronically through Brightspace. It is your responsibility to ensure that your
% proposal is submitted properly. Students may submit the proposal earlier to receive earlier feedback on the proposal.
% The proposal is worth 12% of the project grade.

\section{Background and Objective}

\begin{enumerate}
    \item 
    \item 
    \item 
\end{enumerate}

\section{Methods and Datasets}


\section{Analysis Strategy}


\section{Novelty of Project}


\section{Timeline and Milestones}

\begin{itemize}
    \item \textbf{Week 1 (February 23 - March 1):}
    \item \textbf{Week 2 (March 2 - March 8):}
    \item \textbf{Week 3 (March 9 - March 15):}
    \item \textbf{Week 4 (March 16 - March 22):}
    \item \textbf{Week 5 (March 23 - March 29):}
\end{itemize}